\section{Conclusion}
\label{sec:conclusion}

% CLAIM: C01,C02,C05,C09,C11,C12
A position sample, a future-reference objective, an executable next-cycle
target, and an executed command are not interchangeable target states.  By
assigning each an explicit represented and availability time, the formulation
separates causal estimation, prediction, governance, following, and plant
execution.  The accompanying feasibility taxonomy prevents componentwise
admissibility from being mistaken for one-step reachability or
bounded-jerk sequence consistency.  Under the stated triple-integrator
assumptions, the executable-target governor constructs an adjacent command by
selecting and exactly integrating constant jerk, while profile-aware auditing
preserves the distinct identities of native, shielded, direct, and fallback
execution.

% CLAIM: C03,C04,C06,C07,C08,C10,C14,C17,C19
The controlled evidence supports equally bounded conclusions.  Reliable
analytic target velocity helps the three smooth Phase~A references; analytic
acceleration adds no observed position benefit beyond velocity in those
conditions.  A noncausal next-cycle analytic oracle serves as an indexing
sanity control, while every finite-difference-derived pipeline tested on the
single fixed-grid development trace is worse than the position-only baseline.
Frozen synthetic direct commands record zero audited violations, projections,
and fallbacks, and their separate runtime population records no registered
deadline miss.  A fresh locked V4 attempt produced a large observed RMSE
difference but did not establish a confirmatory PVA benefit because the
preregistered validity and runtime gates were not all satisfied.  None of these
observations ranks a corrected ordinary follower against the direct method or
establishes hardware safety.

% CLAIM: C13,N01,N02,N03
The broader lesson is procedural as well as algorithmic: target information,
executed profile, fallback identity, evidence class, and denominator must be
reported together.  The current work has no independent real-stream
confirmation, hardware evaluation, or dynamics or torque analysis.  V4 is a
frozen non-confirmatory result and may not be rerun on the same test.  A newly
designed confirmation and independent real or hardware evaluation remain
necessary before performance or deployment claims.
